\section{Introduction}
\label{sec:introduction}

Measures of dependence reduce a joint distribution to a small number of
interpretable summaries. Pearson's correlation quantifies linear association,
the rank coefficients of \citet{spearman1904} and \citet{kendall1938} describe
monotone association, and Blomqvist's beta gives a simple quadrant-based
summary \citep{blomqvist1950}. None of these coefficients characterizes
independence. Hoeffding's rank statistic compares the joint distribution with
the product of its marginals and yields a consistent test of independence
under broad conditions \citep{hoeffding1948}. More recently,
\citet{chatterjee2021} introduced a directional rank coefficient that is zero
if and only if the variables are independent and one if and only if the
response is almost surely a measurable function of the predictor.

Many applications attach a nonnegative weight to every observation: survey and
inverse-probability weights, importance weights, or local weights that focus
an analysis on a subpopulation. For means and covariances, weighting amounts
to replacing the empirical distribution by a weighted empirical distribution.
For rank statistics of order two or higher, this substitution is not enough.
Products containing repeated observations must be removed, ties require
compatible weighted counts, and the variability of a weighted statistic
depends on the concentration of the weights rather than only on the nominal
sample size. The effective sample size of \citet{kish1965} is exact for the
conditional variance of a weighted mean, but it is not automatically the
correct variance correction for nonlinear rank statistics.

This report develops a common case-weight framework for six dependence
measures and their independence tests. The methods were implemented several
years ago in the header-only C++ library \texttt{wdm} and its R interface,
but their derivations have not been recorded in writing. Two of them do not
appear to have been studied before, and they are the focus of the report.

The first is a weighted version of Hoeffding's $D$ together with an
$O(n\log n)$ algorithm. The estimator averages over tuples of distinct
observations, each weighted by the product of its case masses, which keeps
it unbiased under exogenous weights. Proposition~\ref{prop:weighted-hoeffding-ranks}
writes its order-three to order-five sums in terms of weighted marginal ranks,
weighted bivariate ranks computed with powers of the masses, and elementary
symmetric sums of the masses. The bivariate ranks are obtained by a weighted
dominance merge sort, so the whole statistic costs a fixed number of sorts
instead of an enumeration of tuples. Under diffuse exogenous weights, its
degenerate null distribution converges to the classical quadratic Gaussian
series after an explicit effective-sample-size correction
(Theorem~\ref{thm:weighted-hoeffding-null}).

The second is a weighted extension of Chatterjee's coefficient. Each case
contributes to the response-rank path in proportion to its mass, and the
coefficient reduces exactly to Chatterjee's statistic under uniform weights.
Conditional on the predictor order and the weights, the null path statistic
has an exact mean and a closed-form first-order variance, and the resulting
studentized statistic is asymptotically standard normal
(Theorem~\ref{thm:weighted-chatterjee-null}). Because the argument conditions
on the predictor, the weights may depend on the predictor.

The Pearson, Spearman, Kendall, and Blomqvist coefficients are included to fix
a common notation and to serve as comparators. Their weighted versions and
$O(n\log n)$ algorithms are routine, and we separate their proved first-order
limits from the tie corrections and finite-sample approximations used in the
software. The software itself is a reproducible implementation of the
methods; its interface is documented in the package.

Section~\ref{sec:framework} fixes notation for case weights and weighted
ranks. Section~\ref{sec:standard-measures} records the four standard
coefficients. Sections~\ref{sec:weighted-hoeffding}
and~\ref{sec:weighted-chatterjee} develop the weighted Hoeffding and
Chatterjee coefficients and their algorithms.
Section~\ref{sec:independence-tests} treats the six independence tests, and
Section~\ref{sec:numerical-experiments} examines finite-sample calibration and
power in a compact simulation study. Proofs are collected in the appendices.
